\section{Objetivo general}

El objetivo del trabajo es diseñar, implementar y verificar un sistema digital
que integre un procesador \tech{RISC-V} de 32 bits con pipeline de cinco etapas,
una unidad de debug y un enlace \tech{UART}. El sistema debe recibir un
programa, cargarlo en la FPGA, ejecutarlo en modo continuo o paso a paso y
permitir que una aplicación de PC observe y controle la ejecución sin volver a
sintetizar el diseño.

\section{Objetivos específicos}

\begin{itemize}
  \item Implementar un subconjunto de \tech{RV32I} de 32 instrucciones más la
    instrucción \code{halt}.
  \item Organizar la ejecución en las etapas \tech{IF}, \tech{ID}, \tech{EX},
    \tech{MEM} y \tech{WB}, con cuatro registros interetapa y avance
    \emph{in-order}.
  \item Resolver dependencias de datos mediante forwarding, bypass y stalls.
  \item Implementar memorias de instrucciones y de datos independientes.
  \item Implementar la carga por \tech{UART}, el control de ejecución y la
    observación de PC, banco de registros, latches, memoria utilizada y
    metadatos de fault.
  \item Proveer un mecanismo para ensamblar código en assembly y una
    aplicación de PC que permitan editar, traducir, cargar, ejecutar
    y depurar programas.
  \item Verificar cada bloque de forma aislada y después en los niveles de
    integración, CPU, sistema y placa que correspondan.
  \item Documentar las decisiones de implementación, las limitaciones y la
    evidencia necesaria para la aceptación final.
\end{itemize}

\section{Criterios de aceptación}

La aceptación final se apoya en los requisitos y en la matriz de
cobertura definidas en la documentación. De forma resumida, la
entrega debe demostrar:

\begin{itemize}
  \item que el procesador ejecuta las instrucciones del subconjunto aprobado y
    conserva el orden arquitectónico;
  \item que los hazards representativos producen el resultado correcto y los
    efectos laterales esperados;
  \item que las memorias respetan capacidad, alineación, rango, cero lógico y
    atomicidad de stores;
  \item que la observación de Debug es coherente y no agrega ciclos
    innecesarios;
  \item que la implementación física cumple los requerimientos de timing;
  \item que la demostración recorre carga, ejecución paso a paso,
    ejecución continua, un hazard, observación y reprogramación.
\end{itemize}

Los criterios detallados y los identificadores de requisito se
encuentran en \doc{requirements.md}{la especificación de requisitos}
y en \doc{verification/coverage.md}{la matriz de cobertura}.
